\chapter*{B2\quad 西华选择题题库（精简版）}
\addcontentsline{toc}{chapter}{B2\quad 西华选择题题库（精简版）}
\markboth{B2\quad 西华选择题题库（精简版）}{}

\begin{tishi}
本库是扫描件《10. 西华选择题题库（精简版）》的文字整理稿，为\textbf{赵琴教材配套选择题库}，
全书 6 章共 \textbf{74 题}，全部为单项选择，是 818 选择/判断类题目的主要来源之一
（E5；E1 slide16"选择/判断考概念辨析"）。

\textbf{章号与赵琴教材的对应关系}（原库章号与教材章号不完全一致，做题时按教材章复习）：

\begin{center}
\begin{tabularx}{\textwidth}{@{}lXl@{}}
\toprule
本库章号 & 本库章名 & 对应赵琴教材 \\
\midrule
第一章 & 绪论 & 第 1 章 绪论 \\
第二章 & 流体静力学 & 第 2 章 流体静力学 \\
第三章 & 流体动力学基础 & 第 3--4 章（运动学 $+$ 动力学基础） \\
第四章 & 管嘴、孔口、管路水力计算 & 第 5 章（含水击、长管与并联管路） \\
第五章 & 相似理论与量纲分析 & 第 6 章 \\
第六章 & 流体动力学（无旋流动） & 第 7 章 理想流体动力学 \\
\bottomrule
\end{tabularx}
\end{center}

\textbf{答案说明}：原库第 6 页为手写答案页，已逐题转录进书末参考答案。
本册题号沿用原库编号，一律写作 1-1、2-5 这种"章号-题号"的形式，
不使用小数点写法，以免与章节号混淆。
\end{tishi}

\section*{第一章 绪论（原库第一章，8 题；对应赵琴第 1 章）}

\begin{ti}

\item[1-1] 在研究流体运动时，按照是否考虑流体的黏性，可将流体分为（\quad）。
\begin{choices}
\item 牛顿流体及非牛顿流体
\item 可压缩流体与不可压缩流体
\item 均质流体与非均质流体
\item 理想流体与实际流体
\end{choices}
\xstarfull{2}\yiju{E5 题库 1-1；E4 2022 年 818 单选第 1--2 题（连续介质假设、不可压缩流体判据）；E1 slide33"第一章绪论＝重点掌握基础概念"}\nandu{易}

\item[1-2] 流体是（\quad）一种物质。
\begin{choices}
\item 不断膨胀直到充满容器的
\item 实际上是不可压缩的
\item 不能承受剪切力的
\item 在任一剪切力的作用下不能保持静止的
\end{choices}
\xstarfull{2}\yiju{E5 题库 1-2；E1 slide33"第一章绪论＝重点掌握基础概念"；E1 slide16"选择/判断考概念辨析"}\nandu{易}

\item[1-3] 流体的切应力（\quad）。
\begin{choices}
\item 当流体处于静止状态时不会产生
\item 当流体处于静止状态时，由于内聚力可以产生
\item 仅仅取决于分子的动量交换
\item 仅仅取决于内聚力
\end{choices}
\xstarfull{2}\yiju{E5 题库 1-3；教材式 1.9、1.10 与"静止流体不承受切应力"的结论（E6）}\nandu{易}

\item[1-4] 水力学中，单位质量力是指作用在单位（\quad）液体上的质量力。
\begin{choices}
\item 面积
\item 体积
\item 质量
\item 重量
\end{choices}
\xstarfull{3}\yiju{E5 题库 1-4 与题库选择题第 1 题（单位质量力）；E4 2015 年试题填空第 1 题（重度与动力黏度换算）；E1 slide33"第一章＝重点掌握基础概念"}\nandu{易}

\item[1-5] 不同的液体其黏滞性\underline{\hspace{2em}}，同一种液体的黏滞性具有随温度\underline{\hspace{2em}}而降低的特性。（\quad）
\begin{choices}
\item 相同；降低
\item 相同；升高
\item 不同；降低
\item 不同；升高
\end{choices}
\xstarfull{3}\yiju{E4 2022 年 818 单选第 2 题（不可压缩流体判据，黏性随温度变化是基本概念）；E5 题库 1-6；E1 slide16"填空考名词定义与小计算"；教材表 1.1 与 1.3.3 节结论（E6）}\nandu{易}

\item[1-6] 液体黏度随温度的升高而\underline{\hspace{2em}}，气体黏度随温度的升高而\underline{\hspace{2em}}。（\quad）
\begin{choices}
\item 减小；升高
\item 增大；减小
\item 减小；不变
\item 减小；减小
\end{choices}
\xstarfull{3}\yiju{E5 题库 1-6；教材 1.3.3 节液体黏性与气体黏性随温度变化的相反趋势（E6）；E1 slide33"第一章绪论＝重点掌握基础概念"}\nandu{易}

\item[1-7] 动力黏滞系数的单位是（\quad）。
\begin{choices}
\item $\mathrm{N\cdot s/m}$
\item $\mathrm{N\cdot s/m^2}$
\item $\mathrm{m^2/s}$
\item $\mathrm{m/s}$
\end{choices}
\xstarfull{3}\yiju{E4 2015 年试题填空第 1 题（动力黏度的量纲与单位换算）；E5 题库 1-7；E1 slide33"第六章＝量纲转化"}\nandu{易}

\item[1-8] 下列说法正确的是（\quad）。
\begin{choices}
\item 液体不能承受拉力，也不能承受压力
\item 液体不能承受拉力，但能承受压力
\item 液体能承受拉力，但不能承受压力
\item 液体能承受拉力，也能承受压力
\end{choices}
\xstarfull{2}\yiju{E5 题库 1-8；教材第 1 章"流体几乎不能承受拉力和抵抗拉伸变形"（E6）；E4 历年真题未单独成题}\nandu{易}

\end{ti}

\section*{第二章 流体静力学（原库第二章，10 题；对应赵琴第 2 章）}

\begin{ti}

\item[2-1] 在重力作用下静止液体中，等压面是水平面的条件是（\quad）。
\begin{choices}
\item 同一种液体
\item 相互连通
\item 不连通
\item 同一种液体，相互连通
\end{choices}
\xstarfull{4}\yiju{E5 题库 2-1；E1 slide33"第二章流体静力学＝计算题基础入门、重点掌握方法"；E6 教材第 2 章等压面判据}\nandu{易}

\item[2-2] 压力表的读值是（\quad）。
\begin{choices}
\item 绝对压强
\item 绝对压强与当地大气压的差值
\item 绝对压强加当地大气压
\item 当地大气压与绝对压强的差值
\end{choices}
\xstarfull{4}\yiju{E5 题库 2-2；E1 slide33"第二章＝计算题基础入门、重点掌握方法"；E6 教材 2.3 节相对压强与真空度}\nandu{易}

\item[2-3] 相对压强是指该点的绝对压强与（\quad）的差值。
\begin{choices}
\item 标准大气压
\item 当地大气压
\item 工程大气压
\item 真空压强
\end{choices}
\xstarfull{4}\yiju{E5 题库 2-3；E1 slide33"第二章＝计算题基础入门、重点掌握方法"；E6 教材 2.3 节（压强基准的选取）}\nandu{易}

\item[2-4] 在等角速度旋转液体中（\quad）。
\begin{choices}
\item 各点的测压管水头等于常数
\item 各点的测压管水头不等于常数，但测压管高度等于常数
\item 各点的压强随水深的变化是线性关系
\item 等压面与质量力不一定正交
\end{choices}
\xstarfull{3}\yiju{E5 题库 2-4；E6 教材第 2 章相对平衡（等角速度旋转容器）；E4 历年真题未直接出现}\nandu{中}

\item[2-5] 图示一半径为 $R$ 的圆柱形容器，内盛密度为 $\rho$ 的液体，若容器以等角速度 $\omega$ 绕 $OZ$ 轴旋转，
则 $A$ 点的压强为（\quad）。
\tuyi{坐标原点 $O$ 取在旋转抛物面自由液面的最低点（位于 $OZ$ 轴上），
$OZ$ 轴竖直向上；容器半径为 $R$、液深为 $h$；$A$ 点在容器底面与侧壁的交点处，即半径 $r=R$、
位置 $z=-h$。自由液面方程为 $z=\dfrac{\omega^2r^2}{2g}$，边缘液面升高 $\dfrac{\omega^2R^2}{2g}$。}
\begin{choices}
\item $p_0+\rho gh$
\item $p_0+\rho gh+\dfrac{1}{2}\rho\omega^2R^2$
\item $p_0+\dfrac{1}{2}\rho\omega^2R^2$
\item 以上都不是
\end{choices}
\xstarfull{4}\yiju{E5 题库 2-5；E1 slide33"第二章＝计算题基础入门、重点掌握方法"；E6 教材第 2 章等角速度旋转容器压强分布式}\nandu{中}

\item[2-6] 对于相对平衡液体，描述正确的是（\quad）。
\begin{choices}
\item 等压面与质量力不正交
\item 等压面不可能为水平面
\item 等压面的形状与液体密度有关
\item 两种液体的交界面为等压面
\end{choices}
\xstarfull{3}\yiju{E5 题库 2-6；E6 教材第 2 章等压面性质（正交性、与密度无关）；E4 历年真题未直接出现}\nandu{中}

\item[2-7] 在均质连通的静止液体中，任一（\quad）上各点压强必然相等。
\begin{choices}
\item 平面
\item 水平面
\item 斜面
\item 以上都不对
\end{choices}
\xstarfull{4}\yiju{E5 题库 2-7；E1 slide33"第二章＝计算题基础入门、重点掌握方法"；E6 教材 2.4 节静水压强分布}\nandu{易}

\item[2-8] 静止液体中静水压强与（\quad）的一次方成正比。
\begin{choices}
\item 淹没深度
\item 位置高度
\item 表面压强
\item 以上都不对
\end{choices}
\xstarfull{4}\yiju{E5 题库 2-8；E1 slide33"第二章＝重点掌握方法"；E6 教材 2.4 节 $p=p_0+\rho gh$}\nandu{易}

\item[2-9] 等压面与质量力的关系是（\quad）。
\begin{choices}
\item 平行
\item 斜交
\item 正交
\item 无关
\end{choices}
\xstarfull{3}\yiju{E5 题库 2-9；E6 教材第 2 章"等压面与质量力正交"；E1 slide33"第二章＝重点掌握方法"}\nandu{易}

\item[2-10] 露天水池，水深 $5\ \mathrm{m}$ 处的相对压强为（\quad）。
\begin{choices}
\item $5\ \mathrm{kPa}$
\item $49\ \mathrm{kPa}$
\item $147\ \mathrm{kPa}$
\item $205\ \mathrm{kPa}$
\end{choices}
\xstarfull{4}\yiju{E5 题库 2-10；E1 slide16"填空考名词定义与小计算"；E1 slide33"第二章＝计算题基础入门、重点掌握方法"（$\gamma_{\text{水}}$ 取 $9.8\ \mathrm{kN/m^3}$）}\nandu{易}

\end{ti}

\section*{第三章 流体动力学基础（原库第三章，21 题；对应赵琴第 3--4 章）}

\begin{ti}

\item[3-1] 欧拉法是研究（\quad）。
\begin{choices}
\item 每个质点的速度
\item 每个质点的运动轨迹
\item 固定空间每一点上流体的运动要素随时间的变化规律
\item 每个空间点的质点轨迹
\end{choices}
\xstarfull{2}\yiju{E5 题库 3-1；E2"第 1--3 章偏概念，掌握基本公式与解释即可"；E6 教材第 3 章欧拉法}\nandu{易}

\item[3-2] 如图所示，等直径管，考虑损失，$A$-$A$ 断面为过水断面，$B$-$B$ 断面为水平面，
试确定 1、2、3、4 各点的物理量有以下关系（\quad）。
\tuyi{等直径倾斜管道，水流由断面 1-1 流向 2-2；$A$-$A$ 为过水断面（与管轴正交），
$B$-$B$ 为水平面；点 1、2 在同一过水断面 $A$-$A$ 上，点 3、4 在同一过水断面 2-2 上，
且点 1、3 在同一水平面 $B$-$B$ 上（故 $z_1=z_3$），点 2、4 为各自断面的形心。}
\begin{choices}
\item $z_1+\dfrac{p_1}{\rho g}=z_2+\dfrac{p_2}{\rho g}$
\item $p_3=p_4$
\item $z_3+\dfrac{p_3}{\rho g}=z_4+\dfrac{p_4}{\rho g}$
\item $p_1=p_2$
\end{choices}
\xstarfull{2}\yiju{E5 题库 3-2；E4 2022 年 818 单选第 6 题（伯努利方程所选计算点位置）；E6 教材第 4 章总流伯努利方程}\nandu{中}

\item[3-3] 恒定流是（\quad）。
\begin{choices}
\item 流动随时间按一定规律变化的流动
\item 流场中任意空间点的运动要素不随时间变化的流动
\item 各过水断面的速度分布相同的流动
\item 各过水断面的压强相同的流动
\end{choices}
\xstarfull{2}\yiju{E5 题库 3-3；E6 教材第 3 章恒定流定义；E2"第 1--3 章偏概念"}\nandu{易}

\item[3-4] 在恒定流条件下（\quad）。
\begin{choices}
\item 流线和迹线正交
\item 流线和迹线重合
\item 流线是平行直线
\item 迹线是平行直线
\end{choices}
\xstarfull{2}\yiju{E5 题库 3-4 与题库选择题第 5 题（流线与迹线重合）；E6 教材第 3 章流线与迹线}\nandu{易}

\item[3-5] 二维流动的速度分布为 $u_x=2x+t$，$u_y=y^2+2t$，则 $t=0$ 时，
点 $(1,2)$ 的液体加速度分量 $a_x$ 和 $a_y$ 分别为（\quad）。
\begin{choices}
\item 1；2
\item 3；16
\item 2；4
\item 5；18
\end{choices}
\xstarfull{2}\yiju{E5 题库 3-5；E6 教材第 3 章质点导数（当地加速度与迁移加速度）；E4 历年真题未直接出现}\nandu{中}

\item[3-6] 连续性方程表示控制体的（\quad）守恒。
\begin{choices}
\item 能量
\item 动量
\item 流量
\item 质量
\end{choices}
\xstarfull{2}\yiju{E5 题库 3-6；E6 教材第 3 章连续性方程；E2"第 1--3 章偏概念"}\nandu{易}

\item[3-7] 变直径管道，直径 $d_1=320\ \mathrm{mm}$，$d_2=160\ \mathrm{mm}$，流速 $v_1=1.5\ \mathrm{m/s}$，
则 $v_2$ 等于（\quad）。
\begin{choices}
\item $3\ \mathrm{m/s}$
\item $4\ \mathrm{m/s}$
\item $6\ \mathrm{m/s}$
\item $9\ \mathrm{m/s}$
\end{choices}
\xstarfull{3}\yiju{E4 2015 年试题填空"变径管 $v_2$"；E5 题库 3-7；E1 slide16"填空考名词定义与小计算"}\nandu{易}

\item[3-8] 通过一个曲面上的体积流量与曲面上的（\quad）有关。
\begin{choices}
\item 法向速度
\item 切向速度
\item 密度分布
\item 压强
\end{choices}
\xstarfull{2}\yiju{E5 题库 3-8；E6 教材第 3 章流量定义 $\dd Q=v_n\,\dd A$；E2"第 1--3 章偏概念"}\nandu{易}

\item[3-9] 伯努利方程中 $z+\dfrac{p}{\rho g}+\dfrac{\alpha v^2}{2g}$ 表示（\quad）。
\begin{choices}
\item 单位重量流体具有的机械能
\item 单位质量流体具有的机械能
\item 单位体积流体具有的机械能
\item 通过过流断面流体的总机械能
\end{choices}
\xstarfull{3}\yiju{E4 2015 年试题填空"伯努利方程各项能量含义"；E5 题库 3-9；E1 slide16"填空考名词定义"}\nandu{易}

\item[3-10] 流量一定，管径沿程减小时，测压管水头线（\quad）。
\begin{choices}
\item 可能沿程上升也可能沿程下降
\item 总是与总水头线平行
\item 只能沿程下降
\item 不可能低于管轴线
\end{choices}
\xstarfull{3}\yiju{E4 2022 年 818 单选第 3 题（测压管水头线坡度小于 0 的沿程变化）；E5 题库 3-10 与题库选择题第 9 题}\nandu{中}

\item[3-11] 黏性流体总水头线沿程的变化是（\quad）。
\begin{choices}
\item 沿程下降
\item 沿程上升
\item 保持水平
\item 前三种情况都有可能
\end{choices}
\xstarfull{3}\yiju{E4 2022 年 818 单选第 3 题；E5 题库 3-11；E6 教材第 4 章水头损失与总水头线}\nandu{易}

\item[3-12] 如图，水平放置的渐扩管，如忽略水头损失，断面形心点的压强有以下关系（\quad）。
\tuyi{水平放置的渐扩管，断面 1-1（小）与 2-2（大）为渐变流过水断面；
点 1、2 分别在两断面的形心处（两断面形心在同一水平线上，故 $z_1=z_2$）；
水流由断面 1-1 流向断面 2-2，因管径扩大而流速减小。}
\begin{choices}
\item $p_1>p_2$
\item $p_1<p_2$
\item $p_1=p_2$
\item 不定
\end{choices}
\xstarfull{3}\yiju{E4 2022 年 818 单选第 6 题（伯努利方程计算点）；E5 题库 3-12；E6 教材第 4 章伯努利方程的能量转换}\nandu{中}

\item[3-13] 在（\quad）流动中，伯努利方程不成立。
\begin{choices}
\item 恒定
\item 理想流体
\item 不可压缩
\item 可压缩
\end{choices}
\xstarfull{3}\yiju{E4 2022 年 818 单选第 6 题；E5 题库 3-13；E6 教材第 4 章伯努利方程的适用条件}\nandu{易}

\item[3-14] 速度水头的表达式为（\quad）。
\begin{choices}
\item $\sqrt{2gh}$
\item $\rho v^2/2$
\item $v^2/(2g)$
\item $v^2/2$
\end{choices}
\xstarfull{2}\yiju{E5 题库 3-14；E4 2015 年试题填空"伯努利方程各项能量含义"；E6 教材第 4 章}\nandu{易}

\item[3-15] 动能修正系数 $\alpha$ 等于（\quad）。
\begin{choices}
\item $\displaystyle\frac{1}{A}\int_A\frac{u}{v}\,\dd A$
\item $\displaystyle\frac{1}{A}\int_A\left(\frac{u}{v}\right)^2\dd A$
\item $\displaystyle\frac{1}{A}\int_A\left(\frac{u}{v}\right)^3\dd A$
\item $\displaystyle\frac{1}{A}\int_A\left(\frac{u}{v}\right)^4\dd A$
\end{choices}
\xstarfull{2}\yiju{E5 题库 3-15；E6 教材第 4 章动能修正系数定义；E4 历年真题未直接出现}\nandu{中}

\item[3-16] 总流的伯努利方程中的速度 $v$ 是（\quad）。
\begin{choices}
\item 某点的速度
\item 断面平均流速
\item 断面形心处的速度
\item 断面上的最大速度
\end{choices}
\xstarfull{2}\yiju{E5 题库 3-16；E6 教材第 4 章总流伯努利方程；E2"第 1--3 章偏概念"}\nandu{易}

\item[3-17] 文丘里管用于测量（\quad）。
\begin{choices}
\item 点速度
\item 压强
\item 密度
\item 流量
\end{choices}
\xstarfull{4}\yiju{E4 2026 年回忆版计算题第 2 题（文丘里管流量证明）；E4 2015 年计算题（文丘里流量计）；E5 题库 3-17}\nandu{易}

\item[3-18] 毕托管用于测量（\quad）。
\begin{choices}
\item 点速度
\item 压强
\item 密度
\item 流量
\end{choices}
\xstarfull{2}\yiju{E5 题库 3-18；E6 教材第 4 章毕托管测速原理；E4 历年真题未直接出现（常与文丘里管并列辨析）}\nandu{易}

\item[3-19] 动量修正系数 $\beta$ 等于（\quad）。
\begin{choices}
\item $\displaystyle\frac{1}{A}\int_A\frac{u}{v}\,\dd A$
\item $\displaystyle\frac{1}{A}\int_A\left(\frac{u}{v}\right)^2\dd A$
\item $\displaystyle\frac{1}{A}\int_A\left(\frac{u}{v}\right)^3\dd A$
\item $\displaystyle\frac{1}{A}\int_A\left(\frac{u}{v}\right)^4\dd A$
\end{choices}
\xstarfull{3}\yiju{E4 2015 年计算题"水平弯管动量方程求 $F_x$、$F_y$"；E5 题库 3-19；E4 2026 年回忆版计算题第 3 题（书上动量方程例题 4.12）}\nandu{中}

\item[3-20] 应用总流的伯努利方程时，两截面之间（\quad）。
\begin{choices}
\item 必须都是急变流
\item 必须都是渐变流
\item 不能出现急变流
\item 可以出现急变流
\end{choices}
\xstarfull{3}\yiju{E4 2022 年 818 单选第 6 题（计算断面选取）；E5 题库 3-20 与题库选择题第 6 题（"之间可以有急变流"）；E6 教材第 4 章适用条件}\nandu{中}

\item[3-21] 流量为 $Q$、速度为 $v$ 的射流冲击一块与流向垂直的平板，则平板受到的冲击力为（\quad）。
\begin{choices}
\item $Qv$
\item $Qv^2$
\item $\rho Qv$
\item $\rho Qv^2$
\end{choices}
\xstarfull{4}\yiju{E4 2015 年计算题"水平弯管动量方程求 $F_x$、$F_y$"；E4 2026 年回忆版计算题第 3 题（书上动量方程例题 4.12）；E5 题库 3-21}\nandu{中}

\end{ti}

\section*{第四章 管嘴、孔口、管路水力计算（原库第四章，20 题；对应赵琴第 5 章）}

\begin{ti}

\item[4-1] 管道中液体的雷诺数与（\quad）无关。
\begin{choices}
\item 温度
\item 管径
\item 流速
\item 管长
\end{choices}
\xstarfull{4}\yiju{E4 2022 年 818 单选第 9 题（两圆管流态判别）；E4 2015 年试题填空"$\Rey$ 判流态"；E5 题库 4-1}\nandu{易}

\item[4-2] 某圆管直径 $d=30\ \mathrm{mm}$，其中液体平均流速为 $20\ \mathrm{cm/s}$，
液体运动黏度为 $0.0114\ \mathrm{cm^2/s}$，则此管中液体流态为（\quad）。
\begin{choices}
\item 层流
\item 层流向紊流过渡
\item 紊流
\item 无法判断
\end{choices}
\xstarfull{4}\yiju{E4 2022 年 818 单选第 9 题（两圆管流态判别）；E4 2015 年试题填空"$\Rey$ 判流态"；E5 题库 4-2}\nandu{易}

\item[4-3] 等直径圆管中紊流的过流断面流速分布是（\quad）。
\begin{choices}
\item 呈抛物线分布
\item 呈对数线分布
\item 呈椭圆曲线分布
\item 呈双曲线分布
\end{choices}
\xstarfull{4}\yiju{E4 2022 年 818 单选第 9--10 题（流态判别与层流损失）；E4 2015 年试题填空"层流断面平均速度＝0.5 倍最大速度"；E5 题库 4-3}\nandu{易}

\item[4-4] 等直径圆管中的层流，其过流断面平均流速是圆管中最大流速的（\quad）。
\begin{choices}
\item 1.0 倍
\item $1/3$ 倍
\item $1/4$ 倍
\item $1/2$ 倍
\end{choices}
\xstarfull{5}\yiju{E4 2015 年试题填空"层流断面平均速度＝0.5 倍最大速度"（原题直接命中）；E4 2022 年 818 单选第 10 题（层流状态管径减半压强损失倍数）；E5 题库 4-4}\nandu{易}

\item[4-5] 圆管中的层流的沿程损失与管中平均流速的（\quad）成正比。
\begin{choices}
\item 一次方
\item 二次方
\item 三次方
\item 四次方
\end{choices}
\xstarfull{5}\yiju{E1 slide33"第五章孔口/管嘴/管路水力计算＝考试计算题重点"；E2"第五章同为主要计算考点"；E5 题库 4-5；E6 教材第 5 章层流沿程损失 $h_f\propto v$}\nandu{易}

\item[4-6] 圆管的水力半径是（\quad）。
\begin{choices}
\item $d/2$
\item $d/3$
\item $d/4$
\item $d/5$
\end{choices}
\xstarfull{5}\yiju{E1 slide33"第五章＝考试计算题重点"；E2"第五章同为主要计算考点"；E5 题库 4-6；E6 教材第 5 章水力半径 $R=A/\chi$}\nandu{易}

\item[4-7] 谢才公式中谢才系数的单位是（\quad）。\textbf{（原资料此题的答案页空白，本册按题面自行判定）}
\begin{choices}
\item 无量纲
\item $\mathrm{m^{1/2}/s}$
\item $\mathrm{m^{3/2}/s}$
\item $\mathrm{m^2/s}$
\end{choices}
\xstarfull{2}\yiju{E5 题库 4-7（原库答案页此处空缺，按题面自行判定）；E6 教材第 5 章谢才公式 $v=C\sqrt{RJ}$；E4 历年真题未直接出现}\nandu{中}

\item[4-8] 判断层流和紊流的临界雷诺数是（\quad）。
\begin{choices}
\item 上临界雷诺数
\item 下临界雷诺数
\item 上下临界雷诺数代数平均
\item 上下临界雷诺数几何平均
\end{choices}
\xstarfull{2}\yiju{E5 题库 4-8；E4 2015 年试题填空"$\Rey$ 判流态"；E6 教材第 5 章上、下临界雷诺数}\nandu{易}

\item[4-9] 圆管紊流光滑区的沿程阻力系数 $\lambda$ 与（\quad）有关。
\begin{choices}
\item 雷诺数 $\Rey$ 有关
\item 管壁相对粗糙 $k_s/d$ 有关
\item $\Rey$ 和 $k_s/d$ 都有关
\item $\Rey$ 和管长有关
\end{choices}
\xstarfull{4}\yiju{E4 2022 年 818 单选第 10 题（层流损失换算）；E1 slide33"第五章＝考试计算题重点"；E5 题库 4-9；E6 教材第 5 章尼古拉兹实验分区}\nandu{中}

\item[4-10] 圆管流的临界雷诺数（\quad）。
\begin{choices}
\item 随管径变化
\item 随流体的密度变化
\item 随流体的黏度变化
\item 不随以上各量变化
\end{choices}
\xstarfull{4}\yiju{E4 2015 年试题填空"$\Rey$ 判流态"；E1 slide33"第五章＝考试计算题重点"；E5 题库 4-10}\nandu{易}

\item[4-11] 在正常工作条件下，作用水头、直径相同时，小孔口出流量（\quad）圆柱形外管嘴流量。
\begin{choices}
\item 大于
\item 小于
\item 等于
\item 不定
\end{choices}
\xstarfull{4}\yiju{E4 2022 年 818 单选第 8 题（圆柱形直角外管嘴正常工作条件 $H_0\le9\ \mathrm{m}$、$l=(3\sim4)d$）；E1 slide33"第五章＝考试计算题重点"；E5 题库 4-11}\nandu{中}

\item[4-12] 长管并联管道各并联管段的（\quad）相等。
\begin{choices}
\item 水头损失
\item 水力坡度
\item 总能量损失
\item 通过的流量
\end{choices}
\xstarfull{5}\yiju{E1 slide33"第五章＝考试计算题重点"；E2"第五章同为主要计算考点"；E5 题库 4-12；E6 教材第 5 章并联管路的水力特征}\nandu{中}

\item[4-13] 长管是（\quad）的管道。
\begin{choices}
\item 管长大于 $100\ \mathrm{m}$
\item 局部损失和流速水头小于总损失的 $5\%$，可忽略不计
\item 沿程损失系数 $\lambda$ 很大
\item 沿程损失大于 $5\ \mathrm{m}$
\end{choices}
\xstarfull{4}\yiju{E1 slide33"第五章＝考试计算题重点"；E2"第五章同为主要计算考点"；E5 题库 4-13 与题库计算题第 2 题（水箱垂直管道出流）}\nandu{易}

\item[4-14] 水从水池经管道流入大气，计算出流量时，应用能量方程应选用计算断面是（\quad）。
\begin{choices}
\item 管道内的任意断面
\item 管道进、出口
\item 水池水面和管道出口
\item 水池水面和管道进口
\end{choices}
\xstarfull{5}\yiju{E4 2022 年 818 单选第 6 题（伯努利方程所选计算点位置，原题直接命中）；E1 slide33"第五章＝考试计算题重点"；E5 题库 4-14}\nandu{中}

\item[4-15] 一条管道将水从高水位水池引入低水位水池，应用能量方程时应选用计算断面是（\quad）。
\begin{choices}
\item 两水池水面
\item 管道进出口
\item 高水池水面和管道出口
\item 管道进口与低水池水面
\end{choices}
\xstarfull{5}\yiju{E4 2022 年 818 单选第 6 题（计算断面选取）；E1 slide33"第五章＝考试计算题重点"；E5 题库 4-15}\nandu{中}

\item[4-16] 虹吸管最高处压强（\quad）大气压。
\begin{choices}
\item $>$
\item $<$
\item $=$
\item $\ge$
\end{choices}
\xstarfull{4}\yiju{E4 2022 年 818 单选第 7 题（虹吸管最高处压强与大气压的关系，原题直接命中）；E5 题库 4-16 与题库选择题第 8 题}\nandu{易}

\item[4-17] 长管是（\quad）的管道。
\begin{choices}
\item 管长大于 $100\ \mathrm{m}$
\item 沿程损失远大于局部损失
\item 沿程损失系数 $\lambda$ 较大
\item 沿程水头损失大于 $5\ \mathrm{m}$
\end{choices}
\xstarfull{4}\yiju{E1 slide33"第五章＝考试计算题重点"；E2"第五章同为主要计算考点"；E5 题库 4-17（与 4-13 重复考同一概念，可见其高频）}\nandu{易}

\item[4-18] （\quad）是并联管路的计算式。
\begin{choices}
\item $Q=Q_1=Q_2$
\item $h_f=h_{f1}+h_{f2}$
\item $h_{f1}=h_{f2}$
\item $v=v_1+v_2$
\end{choices}
\xstarfull{5}\yiju{E1 slide33"第五章＝考试计算题重点"；E2"第五章同为主要计算考点"；E5 题库 4-18（与 4-12 重复考同一概念）；E6 教材第 5 章}\nandu{中}

\item[4-19] 水击是（\quad）。
\begin{choices}
\item 水流对管壁的压力
\item 水压强发生急剧升高和降低的交替变化现象
\item 水射流对阀门的冲击
\item 水拍打水管
\end{choices}
\xstarfull{5}\yiju{E4 2026 年回忆版简答考"水击现象"（原题直接命中）；E1 slide33"第五章＝考试计算题重点"；E5 题库 4-19}\nandu{易}

\item[4-20] 分析水击现象时，必须考虑（\quad）的影响。
\begin{choices}
\item 水深
\item 水的黏性
\item 流体的压缩性和管壁的弹性
\item 雷诺数
\end{choices}
\xstarfull{5}\yiju{E4 2026 年回忆版简答考"水击现象"（直接关联）；E1 slide33"第五章＝考试计算题重点"；E5 题库 4-20}\nandu{中}

\end{ti}

\section*{第五章 相似理论与量纲分析（原库第五章，10 题；对应赵琴第 6 章）}

\begin{ti}

\item[5-1] 进行模型实验，要实现有压管流的动力相似，应满足（\quad）。
\begin{choices}
\item 雷诺准则
\item 弗劳德准则
\item 欧拉准则
\item 柯西准则
\end{choices}
\xstarfull{3}\yiju{E5 题库 5-1 与题库选择题第 12 题（"有压管流应选雷诺准则"）；E1 slide33"第六章相似理论与量纲分析"；E2"第六到第八章只有少量填空/判断/选择"}\nandu{易}

\item[5-2] 进行模型实验，要实现明渠水流的动力相似，应满足（\quad）。
\begin{choices}
\item 雷诺准则
\item 弗劳德准则
\item 欧拉准则
\item 柯西准则
\end{choices}
\xstarfull{3}\yiju{E5 题库 5-2；E1 slide33"第六章相似理论与量纲分析"；E6 教材第 6 章重力相似与弗劳德准则}\nandu{易}

\item[5-3] 压力输水管同种流体的模型实验，已知长度比为 4，则两者的流量比为（\quad）。
\begin{choices}
\item 2
\item 4
\item 8
\item $1/4$
\end{choices}
\xstarfull{2}\yiju{E5 题库 5-3；E6 教材第 6 章雷诺相似（有压管流）的比尺关系；E4 历年真题未直接出现}\nandu{中}

\item[5-4] 明渠水流模型实验，已知长度比为 4，则两者的流量比为（\quad）。
\begin{choices}
\item 16
\item 4
\item 8
\item 32
\end{choices}
\xstarfull{2}\yiju{E5 题库 5-4；E6 教材第 6 章重力相似（明渠）的比尺关系；E4 历年真题未直接出现}\nandu{中}

\item[5-5] 长度 $l$、速度 $v$、重力加速度 $g$ 的无量纲组合是（\quad）。
\begin{choices}
\item $\dfrac{v}{gl}$
\item $\dfrac{lv}{g}$
\item $\dfrac{v^2}{gl}$
\item $\dfrac{gv}{l}$
\end{choices}
\xstarfull{3}\yiju{E5 题库 5-5；E1 slide33"第六章＝量纲转化、相似理论"；E6 教材第 6 章 $\pi$ 定理与弗劳德数}\nandu{中}

\item[5-6] 压强 $p$、密度 $\rho$、长度 $l$、流量 $Q$ 的无量纲组合是（\quad）。
\begin{choices}
\item $\dfrac{pQ^2}{\rho l}$
\item $\dfrac{\rho Q}{pl}$
\item $\dfrac{plQ}{\rho}$
\item $\dfrac{pl^4}{\rho Q^2}$
\end{choices}
\xstarfull{3}\yiju{E5 题库 5-6；E1 slide33"第六章＝量纲转化"；E4 2015 年试题填空（量纲分析相关）}\nandu{中}

\item[5-7] 雷诺数的物理意义是（\quad）。
\begin{choices}
\item 黏滞力与重力之比
\item 惯性力与重力之比
\item 惯性力与黏滞力之比
\item 压力与黏滞力之比
\end{choices}
\xstarfull{3}\yiju{E4 2022 年 818 单选第 9 题（流态判别须用 $\Rey$）；E4 2026 年回忆版简答考"层流与雷诺数关系"；E5 题库 5-7}\nandu{易}

\item[5-8] 弗劳德数的物理意义是（\quad）。
\begin{choices}
\item 黏滞力与重力之比
\item 惯性力与重力之比
\item 惯性力与黏滞力之比
\item 压力与黏滞力之比
\end{choices}
\xstarfull{2}\yiju{E5 题库 5-8；E6 教材第 6 章弗劳德数；E4 历年真题未直接出现}\nandu{易}

\item[5-9] 设模型比尺为 $1:100$，符合重力相似准则，如果模型流量为 $100\ \mathrm{cm^3/s}$，
则原型流量为（\quad）$\mathrm{m^3/s}$。
\begin{choices}
\item 0.01
\item 108
\item 10
\item 10000
\end{choices}
\xstarfull{3}\yiju{E5 题库 5-9；E1 slide33"第六章＝量纲转化、相似理论"；E6 教材第 6 章重力相似流量比尺 $Q_r=\lambda_l^{5/2}$}\nandu{难}

\item[5-10] 运动黏度的量纲是（\quad）。
\begin{choices}
\item $\mathrm{L}/\mathrm{T}^2$
\item $\mathrm{L}/\mathrm{T}^3$
\item $\mathrm{L}^2/\mathrm{T}$
\item $\mathrm{L}/\mathrm{T}$
\end{choices}
\xstarfull{3}\yiju{E5 题库 5-10 与题库选择题第 11 题（运动黏滞系数的量纲）；E4 2015 年试题填空（量纲换算）；E1 slide33"第六章＝量纲转化"}\nandu{易}

\end{ti}

\section*{第六章 流体动力学（原库第六章，5 题；对应赵琴第 7 章无旋流动）}

\begin{ti}

\item[6-1] 无旋流动是（\quad）。
\begin{choices}
\item 流线是直线的流动
\item 迹线是直线的流动
\item 流体微团无旋转的流动
\item 恒定流动
\end{choices}
\xstarfull{3}\yiju{E5 题库 6-1 与题库选择题第 13 题（涡量为零）；E4 2026 年回忆版计算题第 4 题（证明势函数存在性）；E1 slide33"第七章理想流体动力学＝概念为主"}\nandu{易}

\item[6-2] 平面流动具有流函数的条件是（\quad）。
\begin{choices}
\item 无黏性流体
\item 无旋流动
\item 具有速度势函数
\item 满足连续性方程
\end{choices}
\xstarfull{3}\yiju{E5 题库 6-2；E4 2015 年计算题"流函数求速度与流线"；E1 slide33"第七章＝概念为主"}\nandu{中}

\item[6-3] 流体微团具有（\quad）几种基本运动形式。
\begin{choices}
\item 平移和线变形
\item 平移、线变形和角变形
\item 线变形、角变形和转动
\item 平移、线变形、角变形和转动
\end{choices}
\xstarfull{2}\yiju{E5 题库 6-3；E6 教材第 3 章流体微团运动分解；E2"第六到第八章只有少量填空/判断/选择"}\nandu{易}

\item[6-4] 速度势函数存在的条件是（\quad）。
\begin{choices}
\item 无黏性流体
\item 无旋流动
\item 具有速度势函数
\item 满足连续性方程
\end{choices}
\xstarfull{3}\yiju{E4 2026 年回忆版计算题第 4 题（证明势函数存在性，直接命中）；E5 题库 6-4；E1 slide33"第七章＝概念为主"}\nandu{中}

\item[6-5] 不可压缩流体平面流动中，两条流线的流函数的差值，等于（\quad）。
\begin{choices}
\item 通过该两条流线间的单宽体积流量
\item 通过该两条流线间的流量
\item 通过该两条流线间的单宽质量流量
\item 通过该两条流线间的单宽重量流量
\end{choices}
\xstarfull{3}\yiju{E4 2015 年计算题"流函数求速度与流线"；E5 题库 6-5；E1 slide33"第七章＝概念为主"}\nandu{中}

\end{ti}
